% Standalone no-go paper. This is not a revision of poincare_wall.tex.
\documentclass[11pt]{article}
\usepackage[margin=1.02in]{geometry}
\usepackage{amsmath,amssymb,amsthm,mathtools}
\usepackage{booktabs,array,enumitem,microtype,xurl}
\usepackage[colorlinks=true,linkcolor=blue,urlcolor=blue,citecolor=blue]{hyperref}
\hypersetup{
  pdftitle={Exact-Kernel Fourier Families Obstruct Every Penalty Factored Through a Block Map},
  pdfauthor={Lluis Eriksson},
  pdfsubject={A machine-checked no-go theorem and a quantitative repair criterion for a flat lattice gauge block form},
  pdfkeywords={lattice gauge theory, block averaging, Poincare inequality, Fourier mode, postconditioning, Lean 4}
}

\newtheorem{theorem}{Theorem}[section]
\newtheorem{lemma}[theorem]{Lemma}
\newtheorem{proposition}[theorem]{Proposition}
\newtheorem{corollary}[theorem]{Corollary}
\newtheorem{remark}[theorem]{Remark}
\theoremstyle{definition}
\newtheorem{definition}[theorem]{Definition}

\newcommand{\R}{\mathbb{R}}
\newcommand{\C}{\mathbb{C}}
\newcommand{\Z}{\mathbb{Z}}
\newcommand{\CP}{C_{\mathrm P}}
\newcommand{\ip}[2]{\left\langle #1,#2\right\rangle}
\newcommand{\norm}[1]{\left\lVert #1\right\rVert}
\newcommand{\lean}[1]{\texttt{#1}}
\newcommand{\codepath}[1]{\nolinkurl{#1}}

\emergencystretch=3em
\tolerance=1800

\title{Exact-Kernel Fourier Families Obstruct\\
Every Penalty Factored Through a Block Map\\
in a Flat Lattice Gauge Form}
\author{Lluis Eriksson}
\date{August 4, 2026 (version 1.4)}

\begin{document}
\maketitle

\begin{abstract}
Let a fine periodic lattice have side $LN'$ and let $Q_L$ be the
$L^{-d}$-normalized block average of length-$L$ line integrals.  In four
dimensions, the rescaling $Q_L\mapsto LQ_L$ repairs the elementary
constant-field scaling obstruction to a Poincar\'e estimate.  More generally,
for a power rescaling $L^\sigma Q_L$ in dimension $d$, that sector only forces
$\sigma\ge(d-2)/2$; equality is the unique scale-neutral choice, not a
sufficient coercivity theorem.  We prove a stronger factor-through-kernel
no-go: no scalar functional $\Phi_L(Q_LA)$ with $\Phi_L(0)=0$ can yield
full-space coercivity with a constant uniform in the block side.  The
functional may depend on $L$ and may be nonlinear, discontinuous, anisotropic,
or unbounded; no positivity or growth assumption is used.  For every $L\ge2$,
fixed $N'\ge1$, and $N_c\ge2$, we embed within-block Fourier phases in a real
two-plane of the internal coordinate space and construct transverse
one-cochains.  The first mode $A_L$ satisfies
\[
 Q_LA_L=0,\qquad \operatorname{div}A_L=0,
\]
and, for every dimension $d\ge2$,
\[
 \norm{A_L}^2=(LN')^d,\qquad
 \ip{A_L}{K_0A_L}=(LN')^d\lambda_L,
 \qquad
 \lambda_L=4\sin^2\!\left(\frac{\pi}{L}\right)
 \le \frac{4\pi^2}{L^2}.
\]
Thus every factorized penalty vanishes on the witness.  In the linear case,
the Rayleigh quotient of $K_0+(T_LQ_L)^*(T_LQ_L)$ is exactly $\lambda_L$ for
every scale-dependent continuous linear coarse postconditioner $T_L$, with no
bound on $\norm{T_L}$.  Higher frequencies give an orthogonal $2R$-dimensional
real subspace on which the maximal Rayleigh quotient is
$4\sin^2(\pi R/L)$.  Hence the ordered spectrum of the linear operator has at
least $2R$ eigenvalues below $4\pi^2R^2/L^2$; choosing
$R=\lfloor L^\alpha\rfloor$, $0<\alpha<1$, produces a growing low-energy
cluster rather than a single exceptional line.  Every admissible full-space Poincar\'e constant obeys
$\CP(L)\ge1/\lambda_L\ge L^2/(4\pi^2)$, so no constant chosen before $L$
can work.  We also give an exact repair budget.  If the block map is perturbed
to $\widetilde Q_L$, and a positive term $V_L$ is added, the same witness has
Rayleigh quotient
\[
 \lambda_L+\frac{\norm{T_L\widetilde Q_LA_L}^2}{\norm{A_L}^2}
 +\frac{\ip{A_L}{V_LA_L}}{\norm{A_L}^2}.
\]
Uniform coercivity therefore requires a scale-independent positive lift of
this channel; the condition is also sufficient on the witness line, but not
on the full space.  Here ``volume-uniform'' means uniform while $L$ varies at
an arbitrary fixed positive coarse side $N'$; it is not a claim about
$N'\to\infty$ at fixed $L$.

The construction, kernel identities, exact Hodge energy, factorized-penalty
no-go, repair Rayleigh identity, and necessary repair budget are formalized in
Lean~4.  The focused axiom audit reports only \lean{propext},
\lean{Classical.choice}, and \lean{Quot.sound}.  The higher-frequency spectral
multiplicity theorem is proved analytically here and is not claimed as part of
the formalized surface.  The theorem concerns the stated flat full-domain form; it
does not contradict gauge-restricted propagator constructions and makes no
infinite-volume, continuum, or Yang--Mills mass-gap claim.
\end{abstract}

\noindent\textbf{Keywords.}
Lattice gauge theory; block averaging; Poincar\'e inequality; Fourier mode;
postconditioning; formal verification; Lean~4.

\section{Introduction}

Finite-range positive operators admit Combes--Thomas localization estimates
whose constants depend on the lower spectral edge~\cite{CombesThomas}.  In a
block-spin treatment of lattice gauge fields, a natural finite-dimensional
precision is a flat Hodge operator plus a positive block penalty.  The
question addressed here is sharply limited: can any penalty factored through
the original block variable, in the form
\begin{equation}\label{eq:intro-form}
 \ip{A}{K_0A}+\Phi_L(Q_LA),\qquad \Phi_L(0)=0,
\end{equation}
control the full fine one-cochain norm with a constant independent of the
block side $L$?  The familiar linear postconditioned form
$\Phi_L(B)=\norm{T_LB}^2$ is a special case.

The factor $L$ is not arbitrary.  With the line-integral block map used below,
it exactly balances direction-wise constant fields in dimension four.  The
corresponding general-dimensional arithmetic is often compressed into the
phrase ``critical exponent $(d-2)/2$.''  We separate two statements that this
phrase can obscure: exponents at or above that threshold avoid growth forced
by constant fields, while equality is selected only by exact scale neutrality.
Neither statement is sufficient for full-space coercivity.  The block map also
has mean-zero transverse modes in its exact kernel.  Because a map applied
\emph{after} $Q_L$ cannot distinguish zero from zero, no factorized functional
that vanishes at the coarse zero field can lift them.  This includes every
scalar, anisotropic, mixing, or unbounded scale-dependent linear
postconditioner.  The within-block Fourier family makes the obstruction local
to every block and supplies a growing cluster of low modes, not merely one
adversarial vector.

The main result has eight components.
\begin{enumerate}[label=(\roman*)]
\item The constant-sector calculation isolates exactly what the exponent
  $(d-2)/2$ does and does not imply.
\item The witness is defined on the exact finite periodic lattice and works
  for every $L\ge2$, with no parity restriction.
\item Its cancellation under $Q_L$ and under the flat divergence is exact in
  every block and at every positive coarse side $N'$.
\item Its Hodge Rayleigh quotient is the first-cycle eigenvalue
  $4\sin^2(\pi/L)$, hence is $O(L^{-2})$.
\item Frequencies $1\le r\le R$ produce an orthogonal real subspace of
  dimension $2R$ and an explicit min--max bound on the $2R$-th eigenvalue.
\item The no-go is uniform over arbitrary factorized scalar functionals
  $\Phi_L$ satisfying only $\Phi_L(0)=0$; linear postconditioners are a strict
  subclass.
\item The quantifier order is explicit and machine checked: the family $T_L$
  or $\Phi_L$, and one proposed constant $\CP$, are chosen before all block
  sides $L$.
\item An exact repair identity quantifies how much a changed block map or an
  added positive Hessian must contribute on the missed channel.
\end{enumerate}

An earlier square-wave witness gives the weaker even-scale quotient $8/L$.
It already disproves the same gate, but the Fourier construction identifies
the actual low-frequency mechanism and strengthens the quantitative lower
bound on Poincar\'e constants from linear to quadratic growth.

The conclusion is a domain diagnosis, not a general no-go for lattice gauge
renormalization.  Classical constructions use gauge restrictions, quotient
structures, or background-dependent operators.  The present theorem instead
tests the unrestricted flat form in~\eqref{eq:intro-form}.  A repaired theorem
may still hold on a physically justified sector, after changing the
information captured before postconditioning, or after adding a positive
interacting Hessian that lifts the explicit transverse family.

\section{Finite lattice, block map, and quantifiers}

Fix integers $d\ge2$, $L\ge1$, $N'\ge1$, and $N_c\ge2$.  The fine site set is
\[
 \Lambda_{L,N'}=\left(\Z/(LN')\Z\right)^d,
\]
and a positive bond is $(x,\mu)$ with $x\in\Lambda_{L,N'}$ and
$\mu\in\{0,\ldots,d-1\}$.  The real internal coordinate space is
$E_{N_c}=\R^{N_c^2-1}$.  For a one-cochain $A(x,\mu)\in E_{N_c}$, set
\[
 \norm{A}^2=\sum_{x\in\Lambda_{L,N'}}\sum_{\mu=0}^{d-1}
 \norm{A(x,\mu)}^2.
\]

At the trivial background, let $D_1$ be the discrete exterior derivative and
$\delta$ the discrete divergence.  The flat Hodge operator $K_0$ satisfies
\begin{equation}\label{eq:hodge}
 \ip{A}{K_0A}=\norm{D_1A}^2+\norm{\delta A}^2.
\end{equation}
In the formal development this is derived from the adjoint-defined operator;
it is not introduced as an axiom.

For a coarse site $y\in(\Z/N'\Z)^d$, write
\[
 B_y=\{Ly+r:0\le r_\nu<L\text{ for every }\nu\}.
\]
For a coarse positive bond $(y,\mu)$, define the linearized line-integral
average
\begin{equation}\label{eq:Q}
 (Q_LA)(y,\mu)=L^{-d}\sum_{x\in B_y}\sum_{k=0}^{L-1}
 A(x+ke_\mu,\mu).
\end{equation}
For a continuous linear endomorphism $T_L$ of the coarse one-cochain space,
the tested precision is
\begin{equation}\label{eq:H}
 H_{L,T}=K_0+(T_LQ_L)^*(T_LQ_L).
\end{equation}
This class includes scalar rescaling $T_L=s_LI$, direction-dependent weights,
and mixing of coarse sites, directions, or internal coordinates.  No uniform
operator-norm bound on $T_L$ is imposed.

\begin{proposition}[what the constant-sector wall fixes]
\label{prop:critical-exponent}
Let $A_v(x,\mu)=v_\mu$ be a direction-wise constant field and put
$S_v=\sum_\mu\norm{v_\mu}^2>0$.  Then
\begin{equation}\label{eq:constant-norms}
 \norm{A_v}^2=L^d(N')^dS_v,
 \qquad
 \norm{L^\sigma Q_LA_v}^2=L^{2\sigma+2}(N')^dS_v.
\end{equation}
Consequently, a Poincar\'e estimate with a constant independent of unbounded
$L$ can pass this sector only if
\begin{equation}\label{eq:critical-threshold}
 \sigma\ge\frac{d-2}{2}.
\end{equation}
Equality is the unique exponent for which the two norms in
\eqref{eq:constant-norms} have identical scale dependence; larger exponents
overweight the constant sector and also pass this necessary test.
\end{proposition}

\begin{proof}
There are $(LN')^d$ fine sites.  For a constant field, the inner line sum in
\eqref{eq:Q} supplies $L$, while the $L^{-d}$ normalization cancels the $L^d$
source sites in a block, so $Q_LA_v=L v$ on each coarse bond.  This gives
\eqref{eq:constant-norms}.  Substitution in a Poincar\'e estimate and
cancellation of the positive factor $(N')^dS_v$ yields
$L^{d-2-2\sigma}\le\CP$.  Uniformity for arbitrarily large $L$ forces
\eqref{eq:critical-threshold}; equality is exactly the zero exponent in this
last power.
\end{proof}

\begin{remark}[necessary is not sufficient]\label{rem:not-sufficient}
The threshold in Proposition~\ref{prop:critical-exponent} is sufficient only
for preventing the \emph{constant-sector} lower bound from diverging.  Calling
$\sigma=(d-2)/2$ necessary and sufficient for the full gate would silently add
scale neutrality as a hypothesis and would ignore the kernel exhibited below.
\end{remark}

\begin{definition}[critical full-space gate]\label{def:gate}
In $d=4$, fix $N'\ge1$, $N_c\ge2$, and the adjoint-model parameter defining
the flat maps.  The gate asserts that there is a number $\CP>0$ such that, for
every integer $L\ge1$ and every fine one-cochain $A$,
\begin{equation}\label{eq:gate}
 \norm{A}^2\le\CP\left(\ip{A}{K_0A}+\norm{LQ_LA}^2\right).
\end{equation}
\end{definition}

\begin{definition}[postconditioned full-space gate]\label{def:postgate}
Fix the same parameters and an arbitrary family $(T_L)_{L\ge1}$ of continuous
linear coarse-space endomorphisms.  The postconditioned gate asserts that
there is one $\CP>0$, chosen after the family but before the scale, such that
for every $L\ge1$ and every fine one-cochain $A$,
\begin{equation}\label{eq:postgate}
 \norm{A}^2\le\CP\left(\ip{A}{K_0A}+\norm{T_LQ_LA}^2\right).
\end{equation}
The critical scalar gate is the special case $T_L=LI$ in $d=4$.
\end{definition}

\begin{definition}[factorized full-space gate]\label{def:factorgate}
Fix the same parameters and an arbitrary family of scalar functionals
\[
 \Phi_L:\{\text{coarse one-cochains at side }N'\}\longrightarrow\R,
 \qquad \Phi_L(0)=0.
\]
The factorized gate asserts that there is one $\CP>0$, chosen after the family
but before the scale, such that for every $L\ge1$ and every fine one-cochain
$A$,
\begin{equation}\label{eq:factorgate}
 \norm{A}^2\le\CP\left(\ip{A}{K_0A}+\Phi_L(Q_LA)\right).
\end{equation}
No linearity, continuity, positivity, or growth hypothesis on $\Phi_L$ is
needed for the no-go below.  If ``penalty'' is reserved for nonnegative
functionals, that narrower class is of course included.  Definition
\ref{def:postgate} is recovered by $\Phi_L(B)=\norm{T_LB}^2$.
\end{definition}

\begin{remark}[meaning of volume-uniform]\label{rem:volume}
The coarse side $N'$ is arbitrary but fixed inside one instance of the gate.
The family $(T_L)$ or $(\Phi_L)$ and $\CP$ may depend on $N'$, $N_c$, and the
fixed adjoint model, but not on the subsequently quantified $L$.  The fine volume $(LN')^4$
then varies and the same $\CP$ must work for all of it.  This paper does not call a bound
``volume-uniform'' merely because it is uniform as $N'\to\infty$ at fixed
$L$; that is a different quantifier statement.
\end{remark}

\section{The block-periodic Fourier witness}

Let
\[
 \zeta_L=\exp(2\pi i/L),\qquad L\ge2.
\]
Then $\zeta_L^L=1$, $\zeta_L\ne1$, and the finite geometric sum gives
\begin{equation}\label{eq:rootsum}
 \sum_{r=0}^{L-1}\zeta_L^r=0.
\end{equation}
Choose orthonormal vectors $u_0,u_1\in E_{N_c}$, possible because
$N_c^2-1\ge3$, and define the real-linear isometry
\[
 \iota:\C\longrightarrow E_{N_c},\qquad
 \iota(a+ib)=au_0+bu_1.
\]
For $m\in\Z/(LN')\Z$, define the block-periodic internal profile
\begin{equation}\label{eq:profile}
 \phi_L(m)=\iota\!\left(\zeta_L^{\,m\bmod L}\right).
\end{equation}
It has unit norm and exact zero mean in every length-$L$ block.

Choose distinct directions $i\ne j$ and set
\begin{equation}\label{eq:witness}
 A_L(x,\mu)=
 \begin{cases}
 \phi_L(x_j),&\mu=i,\\
 0,&\mu\ne i.
 \end{cases}
\end{equation}
The field points in direction $i$ and varies only in the transverse direction
$j$.

\begin{lemma}[norm and block kernel]\label{lem:norm-q}
For every $L\ge2$ and $N'\ge1$,
\begin{equation}\label{eq:norm-q}
 \norm{A_L}^2=(LN')^d,
 \qquad Q_LA_L=0.
\end{equation}
Consequently $(T_LQ_L)A_L=0$ for every linear postconditioner $T_L$.
More generally, $\Phi_L(Q_LA_L)=0$ for every scalar functional with
$\Phi_L(0)=0$.
\end{lemma}

\begin{proof}
Only direction-$i$ bonds contribute to the norm, one at each fine site, and
$\norm{\phi_L(m)}=1$.  For the block map, a coarse bond outside direction $i$
sees only zero components.  In direction $i$, line shifts preserve $x_j$.
After the line sum is factored out, the sum over the within-block $j$-offset
is~\eqref{eq:rootsum}.  This cancellation occurs separately in every source
block, hence for every $N'$.
\end{proof}

\begin{lemma}[transverse divergence]\label{lem:div}
At the trivial background,
\[
 \delta A_L=0.
\]
\end{lemma}

\begin{proof}
Only the direction-$i$ component enters.  That component depends only on
$x_j$ and is invariant under forward and backward shifts in direction $i$.
\end{proof}

The witness therefore lies simultaneously in the exact kernel of the block
map and the flat gauge-constraint kernel.  Its remaining energy is purely
curvature energy.

\section{Exact first-mode energy}

Define
\begin{equation}\label{eq:lambda}
 \lambda_L=\lvert\zeta_L-1\rvert^2.
\end{equation}
For every $r$, multiplication by the unit complex number $\zeta_L^r$ gives
\begin{equation}\label{eq:jump}
 \norm{\iota(\zeta_L^r)-\iota(\zeta_L^{r+1})}^2
 =\lvert1-\zeta_L\rvert^2=\lambda_L,
\end{equation}
including the periodic jump $r=L-1$ because $\zeta_L^L=1$.  Elementary
trigonometry yields
\begin{equation}\label{eq:sin}
 \lambda_L=2-2\cos(2\pi/L)
 =4\sin^2(\pi/L)\le\frac{4\pi^2}{L^2}.
\end{equation}

\begin{proposition}[exact flat Hodge energy]\label{prop:energy}
For the transverse witness~\eqref{eq:witness},
\begin{equation}\label{eq:energy}
 \ip{A_L}{K_0A_L}=(LN')^d\lambda_L.
\end{equation}
\end{proposition}

\begin{proof}
Lemma~\ref{lem:div} removes the divergence term in~\eqref{eq:hodge}.  In the
plaquette curl, every direction pair other than the pair containing $i$ and
$j$ vanishes: either the field component is zero or its profile is unchanged
by the shift.  Each $(i,j)$ plaquette contributes the common value
$\lambda_L$ from~\eqref{eq:jump}.  There is one contributing geometric
plaquette at each of the $(LN')^d$ fine sites under the convention used by
$D_1$.  Summation gives~\eqref{eq:energy}.
\end{proof}

\begin{theorem}[postconditioner-invariant Rayleigh quotient]\label{thm:rayleigh}
For every continuous linear coarse-space endomorphism $T_L$,
\begin{equation}\label{eq:rayleigh}
 \frac{\ip{A_L}{H_{L,T}A_L}}{\norm{A_L}^2}
 =\lambda_L=4\sin^2(\pi/L).
\end{equation}
Thus the best full-space coercivity constant $c_L$ of $H_{L,T}$ obeys
\begin{equation}\label{eq:coerc-upper}
 c_L\le\lambda_L\le\frac{4\pi^2}{L^2}.
\end{equation}
\end{theorem}

\begin{proof}
Linearity gives $T_L0=0$, so Lemma~\ref{lem:norm-q} removes the block term
from~\eqref{eq:H}, irrespective of $\norm{T_L}$.  Divide
Proposition~\ref{prop:energy} by the positive norm in
Lemma~\ref{lem:norm-q}; the entire volume factor cancels.
\end{proof}

\section{A spectral ladder, not an isolated witness}

The same cancellation is available at every nonzero within-block frequency.
For $1\le r<L$ and $c\in\C$, define the real one-cochain
\begin{equation}\label{eq:higher-mode}
 A_{r,c}(x,\mu)=
 \begin{cases}
  \iota\!\left(c\exp(2\pi i r x_j/L)\right),&\mu=i,\\
  0,&\mu\ne i,
 \end{cases}
\end{equation}
and let $E_r=\{A_{r,c}:c\in\C\}$.  Each $E_r$ is a real two-dimensional
subspace.  The root-of-unity sum, transverse-divergence calculation, and
plaquette calculation used above give
\begin{equation}\label{eq:higher-identities}
 Q_LA_{r,c}=0,\qquad \delta A_{r,c}=0,\qquad
 \ip{A_{r,c}}{K_0A_{r,c}}=\lambda_{r,L}\norm{A_{r,c}}^2,
 \quad
 \lambda_{r,L}=4\sin^2(\pi r/L).
\end{equation}

\begin{theorem}[growing low-energy multiplicity]\label{thm:ladder}
Let $L\ge3$ and
$1\le R\le\lfloor(L-1)/2\rfloor$.  For every continuous linear
postconditioner $T_L$, the spaces $E_1,\ldots,E_R$ are mutually orthogonal,
and
\[
 S_R:=E_1\oplus\cdots\oplus E_R\subseteq\ker Q_L,
 \qquad \dim_{\R}S_R=2R.
\]
Every nonzero $A\in S_R$ satisfies
\begin{equation}\label{eq:subspace-rayleigh}
 \frac{\ip{A}{H_{L,T}A}}{\norm{A}^2}
 \le\lambda_{R,L}
 =4\sin^2(\pi R/L)
 \le\frac{4\pi^2R^2}{L^2}.
\end{equation}
Consequently, if
$\mu_1(H_{L,T})\le\mu_2(H_{L,T})\le\cdots$ are the eigenvalues counted with
multiplicity, then
\begin{equation}\label{eq:minmax}
 \mu_{2R}(H_{L,T})\le\frac{4\pi^2R^2}{L^2}.
\end{equation}
In particular, for every $0<\alpha<1$ and all sufficiently large $L$, the
choice $R=\lfloor L^\alpha\rfloor$ gives at least
$2\lfloor L^\alpha\rfloor$ eigenvalues bounded by
$4\pi^2L^{-2(1-\alpha)}$.
\end{theorem}

\begin{proof}
The internal embedding satisfies
$\ip{\iota(z)}{\iota(w)}=\operatorname{Re}(z\overline w)$.  Summing the inner
product of an $r$-profile and an $s$-profile over $x_j$ therefore reduces, for
$r\ne s$, to the real part of a nontrivial root-of-unity sum and vanishes.
This proves orthogonality and the dimension count.  Equations
\eqref{eq:higher-identities} show that the block term vanishes and that the
quadratic form is diagonal on the orthogonal sum.  Since
$r\mapsto\sin(\pi r/L)$ is increasing on the stated range, the Rayleigh
quotient of any vector in $S_R$ is at most $\lambda_{R,L}$.  The last
inequality in~\eqref{eq:subspace-rayleigh} is $\sin t\le t$.  The finite
dimensional min--max principle applied to the $2R$-dimensional test subspace
gives~\eqref{eq:minmax}.  Finally,
$\lfloor L^\alpha\rfloor\le(L-1)/2$ for all sufficiently large $L$, and
$R^2/L^2\le L^{-2(1-\alpha)}$.
\end{proof}

Theorem~\ref{thm:ladder} concerns the linear self-adjoint operator
$H_{L,T}$.  For a nonlinear factorized functional $\Phi_L$, an eigenvalue
statement need not make sense; nevertheless, every vector in every $E_r$ is
still erased by $Q_L$, so the domination inequality fails already on $E_1$.
The multiplicity theorem is analytic and is not part of the Lean surface
claimed below.

\section{No uniform critical coercivity}

\begin{theorem}[factor-through-kernel all-scales no-go]\label{thm:nogo}
Fix $N'\ge1$, $N_c\ge2$, the adjoint-model parameter defining the flat
operators, and any family $(\Phi_L)_{L\ge1}$ of scalar functionals on the
coarse field space with $\Phi_L(0)=0$.  The factorized full-space gate in
Definition~\ref{def:factorgate} is false.  More quantitatively, at every
$L\ge2$ any constant valid at that scale must satisfy
\begin{equation}\label{eq:CP-lower}
 \CP(L)\ge\frac{1}{\lambda_L}
 \ge\frac{L^2}{4\pi^2}.
\end{equation}
\end{theorem}

\begin{proof}
Insert $A_L$ into~\eqref{eq:factorgate}.  By Lemma~\ref{lem:norm-q} and
Proposition~\ref{prop:energy},
\[
 (LN')^4\le\CP(LN')^4\lambda_L.
\]
The volume factor is positive, so $1\le\CP\lambda_L$, proving the first
inequality in~\eqref{eq:CP-lower}; the second follows from~\eqref{eq:sin}.
If one $\CP$ worked for all $L$, it would force $L^2\le4\pi^2\CP$ for every
$L$, a contradiction.
\end{proof}

\begin{corollary}[arbitrary linear postconditioning]\label{cor:post}
The postconditioned gate in Definition~\ref{def:postgate} is false for every
family of continuous linear maps $T_L$, even if their norms diverge with $L$.
\end{corollary}

\begin{proof}
Apply Theorem~\ref{thm:nogo} to $\Phi_L(B)=\norm{T_LB}^2$.
\end{proof}

\begin{corollary}[critical and supercritical scalar normalizations]\label{cor:scalar}
No real scalar sequence $s_L$ can restore full-space uniform coercivity for
$K_0+(s_LQ_L)^*(s_LQ_L)$.  In particular, neither the critical choice
$s_L=L^{(d-2)/2}$ nor any faster scalar growth repairs the full-space gate.
\end{corollary}

\begin{proof}
Take $T_L=s_LI$ in Corollary~\ref{cor:post}.
\end{proof}

The point is stronger than ``the critical exponent is insufficient.''  Even
an arbitrary rule depending only on the retained coarse variable acts after
the witness has already been erased.  The proof uses only $\Phi_L(0)=0$;
linearity was never the essential hypothesis.

The theorem isolates the lower spectral edge as the obstruction.  Separate
finite-range, kernel-amplitude, ball-count, and block-metric localization
estimates are not invalidated.  What fails is the positive scale-independent
coercivity input required to turn those estimates into a uniform inverse or
inverse-square-root bound on the full domain.

\section{Exact repair budget on the missed channel}

The kernel theorem says precisely what postconditioning cannot do.  It also
allows a quantitative statement about changes made \emph{before} that kernel
or outside the block penalty.  Let $\widetilde Q_L$ be any linear replacement
for $Q_L$, let $T_L$ be any continuous linear postconditioner, and let $V_L$
be a positive semidefinite self-adjoint operator on fine one-cochains.  Set
\begin{equation}\label{eq:repaired-H}
 \widetilde H_L=K_0+(T_L\widetilde Q_L)^*(T_L\widetilde Q_L)+V_L
\end{equation}
and define the two dimensionless lifts
\begin{equation}\label{eq:lifts}
 \varepsilon_L=
 \frac{\norm{T_L\widetilde Q_LA_L}}{\norm{A_L}},
 \qquad
 \nu_L=\frac{\ip{A_L}{V_LA_L}}{\norm{A_L}^2}\ge0.
\end{equation}

\begin{theorem}[sharp witness-channel repair criterion]\label{thm:repair}
For every $L\ge2$,
\begin{equation}\label{eq:repair-identity}
 \frac{\ip{A_L}{\widetilde H_LA_L}}{\norm{A_L}^2}
 =\lambda_L+\varepsilon_L^2+\nu_L.
\end{equation}
Consequently, if one constant $\CP>0$ gives full-space coercivity for
$\widetilde H_L$ at every scale, then
\begin{equation}\label{eq:repair-necessary}
 \varepsilon_L^2+\nu_L\ge \CP^{-1}-\lambda_L
 \quad(L\ge2),
 \qquad
 \liminf_{L\to\infty}(\varepsilon_L^2+\nu_L)\ge\CP^{-1}.
\end{equation}
Conversely, uniform coercivity restricted to the one-dimensional spaces
$\operatorname{span}\{A_L\}$ holds if and only if
\begin{equation}\label{eq:repair-iff}
 \inf_{L\ge2}(\lambda_L+\varepsilon_L^2+\nu_L)>0.
\end{equation}
This equivalence is only a witness-sector statement; it is not sufficient for
coercivity on the full one-cochain space.
\end{theorem}

\begin{proof}
Expand the three quadratic terms in~\eqref{eq:repaired-H}, use
Proposition~\ref{prop:energy}, and divide by $\norm{A_L}^2$ to obtain
\eqref{eq:repair-identity}.  Testing a full-space Poincar\'e estimate on
$A_L$ gives $1\le\CP(\lambda_L+\varepsilon_L^2+\nu_L)$, hence
\eqref{eq:repair-necessary} because $\lambda_L\to0$.  On the line spanned by
$A_L$, homogeneity makes the displayed Rayleigh quotient the unique quotient,
so a uniform lower bound is equivalent to~\eqref{eq:repair-iff}.
\end{proof}

The exact identity~\eqref{eq:repair-identity} and the finite-scale necessary
budget obtained by testing the full-space estimate are machine checked in
Lean.  The asymptotic liminf statement and the witness-line equivalence
\eqref{eq:repair-iff} are elementary analytic consequences; no full-space
sufficiency claim is made or smuggled into the hypotheses.

The formula distinguishes three logically different repairs.  A
postconditioner after the original $Q_L$ has $\varepsilon_L=0$ exactly and is
therefore powerless.  A new pre-kernel measurement may give
$\varepsilon_L>0$.  An additional Hessian contributes $\nu_L$.  Any proposed
interacting continuation must prove, rather than assume, that their combined
contribution stays uniformly positive on this explicit family.

There is a useful scaling corollary.  Write
$\widetilde Q_L=Q_L+R_L$, take the constant-sector critical scalar
$T_L=L^{(d-2)/2}I$, and put
\begin{equation}\label{eq:leakage}
 \eta_L=\frac{\norm{R_LA_L}}{\norm{A_L}}.
\end{equation}
Because $Q_LA_L=0$, the block contribution is exactly
$\varepsilon_L^2=L^{d-2}\eta_L^2$.  In four dimensions, with $V_L=0$, a
uniform repair of even this one channel requires
\begin{equation}\label{eq:d4-leakage}
 \liminf_{L\to\infty}L^2\eta_L^2>0;
\end{equation}
leakage $o(L^{-1})$ cannot work.  This exposes the scale at which a changed
block observable must see the missing Fourier mode.  It does not manufacture
such an observable or establish its gauge covariance.

\section{Relation to prior gauge-RG constructions}

Ba\l aban's propagator papers study the Gaussian vector-field theory underlying
lattice gauge renormalization and then extend the analysis to operators with
restrictions at several scales and to iterated renormalization
transformations~\cite{BalabanI,BalabanII}.  His separate treatment of averaging
operations studies their regularity on selected classes of gauge fields and
gauge transformations~\cite{BalabanAveraging}.  In that program, block
averages are not used in isolation: gauge conditions, restricted domains,
minimizers, and background-dependent propagators are part of the construction.
Modern expositions of covariant axial gauge likewise emphasize an axial
restriction adapted to block averaging~\cite{DimockAxial}.

There is therefore no conflict between classical positivity statements on a
gauge-fixed or otherwise restricted subspace and Theorem~\ref{thm:nogo}.  The
claims have different domains.  The exact delta established here is:
\begin{quote}
For the explicit finite periodic line-integral map~\eqref{eq:Q}, the flat
full-space form has block-periodic transverse Fourier families in the
simultaneous kernels of $Q_L$ and the flat divergence.  Therefore every term
$\Phi_L(Q_LA)$ with $\Phi_L(0)=0$ misses them exactly.  In the linear case,
$K_0+(T_LQ_L)^*(T_LQ_L)$ has at least $2R$ eigenvalues no larger than
$4\pi^2R^2/L^2$ for $1\le R\le\lfloor(L-1)/2\rfloor$, at every $N'\ge1$.
\end{quote}
This is a no-go for that unrestricted coercivity proposal, not a priority
claim over the general observation that gauge fields require gauge fixing,
and not a no-go for Ba\l aban's restricted propagators.  Indeed, the result
explains why a domain restriction or quotient is mathematically load-bearing
rather than cosmetic.

\section{Machine-checked realization}

The Lean~4 development is an eight-module theorem chain, plus three focused
oracle files, headed by
\begin{center}
\codepath{YangMills/RG/PhysicalCriticalRescalingFourierFactorizedRepair.lean}.
\end{center}
Table~\ref{tab:map} maps the analytic statements to selected declarations.

The constant-sector norm identity and lower bound used in
Proposition~\ref{prop:critical-exponent} are formalized in
\codepath{PhysicalPoincareCriticalRescaling.lean}; the specialization to a
general real exponent $\sigma$ is the elementary power comparison displayed
in the proof.  The Fourier chain then tests the sector that the constant-field
calculation cannot see.

\begin{table}[ht]
\centering
\small
\begin{tabular}{@{}>{\raggedright\arraybackslash}p{0.35\textwidth}
                  >{\raggedright\arraybackslash}p{0.55\textwidth}@{}}
\toprule
Mathematical content & Lean declaration \\
\midrule
Root-of-unity block cancellation &
\codepath{sum_blockFourierRoot_pow_eq_zero} \\
Lie-valued cancellation in every block &
\codepath{sum_blockFourierProfile_block_eq_zero} \\
Eigenvalue bound $\lambda_L\le(2\pi/L)^2$ &
\codepath{blockFourierEigenvalue_le} \\
Exact block-map kernel &
\codepath{flatBlockConstraintQCLM_blockFourierMode_eq_zero} \\
Zero transverse divergence &
\codepath{gaugeConstraintQCLM_blockFourierModeCochain_eq_zero_of_ne} \\
Exact Hodge quadratic form &
\codepath{flatGaugeHodgeK0_inner_blockFourierModeCochain_of_ne} \\
Exact Rayleigh quotient &
\codepath{blockFourierMode_rayleigh_eq} \\
Necessary bound $1\le\CP\lambda_L$ &
\codepath{criticalRescaledFlatPoincare_fourier_lower_bound} \\
All-scales no-go &
\codepath{volumeUniformCriticalRescaledFlatPoincareGate_fourier_false} \\
Arbitrary-postconditioner lower bound &
\codepath{postconditionedFlatPoincare_fourier_lower_bound} \\
Postconditioned all-scales no-go &
\codepath{volumeUniformPostconditionedFlatPoincareGate_fourier_false} \\
Arbitrary factorized-penalty lower bound &
\codepath{factorizedFlatPoincare_fourier_lower_bound} \\
Factor-through-kernel all-scales no-go &
\codepath{volumeUniformFactorizedFlatPoincareGate_fourier_false} \\
Exact repaired witness quotient &
\codepath{repairedFlatPoincare_blockFourier_rayleigh_eq} \\
Necessary repaired full-space budget &
\codepath{repairedFlatPoincare_blockFourier_budget} \\
\bottomrule
\end{tabular}
\caption{Theorem-to-artifact map.}
\label{tab:map}
\end{table}

The formal result uses the exact norm definition
$\lambda_L=\norm{\zeta_L-1}^2$ and proves the bound
$\lambda_L\le(2\pi/L)^2$.  Equation~\eqref{eq:sin} is the elementary
closed-form identification of that norm.  The original focused oracle
\begin{center}
\codepath{YangMills/RG/PhysicalCriticalRescalingFourierNoGoOracle.lean}
\end{center}
reports exactly
\[
 [\lean{propext},\ \lean{Classical.choice},\ \lean{Quot.sound}]
\]
for all nine audited declarations, with no project axiom and no
\lean{sorryAx}.  A second focused oracle,
\begin{center}
\codepath{YangMills/RG/PhysicalCriticalRescalingFourierPostconditionedNoGoOracle.lean},
\end{center}
reports the same list for the block-kernel theorem, the exact Hodge energy,
the eigenvalue estimate, and both postconditioned theorems.  The third oracle,
\begin{center}
\codepath{YangMills/RG/PhysicalCriticalRescalingFourierFactorizedRepairOracle.lean},
\end{center}
reports the same list for six declarations: the reused exact kernel and Hodge
energy, the factorized lower bound and all-scales no-go, and the repaired
Rayleigh identity and necessary budget.

\subsection*{Immutable source and reproduction}

The original seven-file theorem source is fixed at repository commit
\begin{center}
\lean{f21539ed0bb880a04078de369bf5cbf063f7b101}.
\end{center}
Its former head theorem file has LF-normalized SHA-256
\begin{center}
\lean{5F0890D14EA6981CBE6459605A634386ED2C844F6D03A6947E33B34250E2F5AE}.
\end{center}
The v1.4 factorized-and-repair extension has 12,064 LF bytes and SHA-256
\begin{center}
\lean{D7586D24D7B7E9050001AC5A8E4A8061A3EE3F7D441C5373186EBAEB7B97DE38},
\end{center}
and its 537-byte focused oracle has SHA-256
\begin{center}
\lean{6334EE25A6A0D8778A91CC49269BC92563C327DD1BD0E5CA9CF477F5E90738E1}.
\end{center}
The companion packet
\codepath{critical_rescaling_fourier_no_go_sources_v1.4.zip} has SHA-256
\begin{center}
\lean{EF5501145398B5F672DCDFB66206C1F5F280F2C61B231D2C9AC0E1A4E693F82A}.
\end{center}
It contains the eight theorem modules, all three focused oracles, the v1.4
manifest and audit record, the Lean toolchain, and the pinned Lake files.  The
toolchain is Lean \lean{v4.29.0-rc6}.  The pinned Mathlib revision is
\begin{center}
\lean{07642720480157414db592fa85b626dafb71355b}.
\end{center}
From the repository root, run
\begin{quote}
\footnotesize
\lean{lake build}\\
\codepath{YangMills.RG.PhysicalCriticalRescalingFourierFactorizedRepair}
\\[0.3em]
\lean{lake env lean}\\
\codepath{YangMills/RG/PhysicalCriticalRescalingFourierFactorizedRepairOracle.lean}
\end{quote}
The v1.4 reproduction on August 4, 2026 used a Colab Pro+ CPU/high-RAM runtime
(50.99~GB visible RAM) with no GPU, under the visibly confirmed account
\lean{lluiseriksson@gmail.com}.  A disconnected preparation reconstructed the
public base \lean{04f87347f3e4d46a05e77bc1c70855794e111477} plus 17
LF-normalized frozen bodies, checking every body hash before execution.  Lake
update took 124.412 seconds and cache retrieval 6.507 seconds.  The first
target pass materialized the 8,185 imported jobs and exposed a target-local
arithmetic proof defect; the judge and criterion were not changed.  After two
source-only proof repairs, the final 12,064-byte body was frozen at the hash
above.  The final target invocation succeeded with 8,186 jobs in 18.835
seconds and materialized the target \lean{.olean}.  The focused oracle then
exited successfully in 14.953 seconds and returned exactly
$[\lean{propext},\lean{Classical.choice},\lean{Quot.sound}]$ for all six
declarations.  The final frozen run spanned
19:07:36.898702--19:08:10.693578 UTC; the full reconstruction audit spanned
18:56:52.878156--19:08:10.693578 UTC.  The runtime was visibly open no later
than 18:48:10.965571 UTC and was disconnected and deleted immediately after
the audit.  The only two warnings were pre-existing style lints in imported
files.  A local code-token scan found no \lean{sorry}, \lean{admit}, or project
\lean{axiom} in the new source.  The notebook is archived at
\href{https://colab.research.google.com/drive/1iXBNBeBEGr7O3pEbuRtac1VE1mMHX0Lh}
{Colab audit notebook 1iXBNBeBEGr7O3pEbuRtac1VE1mMHX0Lh}.  This is
reproduction evidence and a frozen object for external audit, not an external
certification.

\section{What remains open}

\paragraph{Restricted or quotient sector.}
A repaired domain may remove or quotient the obstructive transverse family,
but the restriction must be physically justified and stable under the
renormalization operations.  Coercivity on that new domain remains to be
proved.

\paragraph{Additional positive Hessian.}
An interacting Wilson Hessian or another positive term $V_L$ could lift the
family.  Theorem~\ref{thm:repair} gives its exact required contribution on
this channel, but establishing a lower bound for the physical interacting
Hessian remains open.  Such a bound is not a consequence of the flat
calculation and is not assumed in the no-go theorem.

\paragraph{Changed block observable.}
A map $\widetilde Q_L$ that measures information discarded by $Q_L$ may lift
the witness.  Equation~\eqref{eq:d4-leakage} gives the necessary four-dimensional
leakage scale for the critical scalar normalization.  Constructing a local,
gauge-covariant, renormalization-compatible map with this property and proving
full-space coercivity remain open.

\paragraph{No continuum conclusion.}
The theorem is finite-dimensional.  It proves no positive coercivity for an
interacting theory, no infinite-volume measure, no continuum limit, no
Osterwalder--Schrader reconstruction, and no Yang--Mills mass gap.

\section{Conclusion}

Critical scalar rescaling can correct a normalization mismatch but cannot
remove an exact kernel; neither can any functional of the block variable
performed after that kernel has formed, provided it vanishes at the coarse
zero field.  The constant-sector wall selects
$\sigma=(d-2)/2$ only when exact scale neutrality is required; uniformity of
that sector alone allows every larger exponent.  None of these scalar choices
acts on $\ker Q_L$.  The first within-block Fourier phase produces, at
every $L\ge2$ and every positive coarse side, a transverse divergence-free
cochain annihilated by $Q_L$.  Its exact Rayleigh quotient is
$4\sin^2(\pi/L)$, so the best full-space spectral lower bound is at most
$4\pi^2/L^2$.  Higher frequencies give a $2R$-dimensional orthogonal test
space below $4\pi^2R^2/L^2$, hence a growing low-energy spectral cluster in
the linear problem.  The current flat full-domain gate is therefore
impossible for every factorized penalty $\Phi_L(Q_LA)$ with $\Phi_L(0)=0$;
continuity, linearity, positivity, and bounded growth play no role.

The negative result is operational: it identifies the precise input that
must change.  A successful continuation must justify a different domain,
measure the missing channel before postconditioning, or add an operator that
supplies uniform energy to the explicitly exhibited Fourier sector.  The
machine-checked repair identity and necessary budget quantify this last
requirement without assuming full-space sufficiency.

\appendix
\section{The square-wave precursor}

For even $L=2M$, the real profile that equals $+1$ on the first half of each
block and $-1$ on the second half also has zero block mean.  Its two jumps per
block have squared magnitude four, giving total one-dimensional energy
$8N'$ and Rayleigh quotient $8/L$.  This witness is useful because it requires
only a single internal vector and elementary interface counting.  The Fourier
witness is stronger in three ways: it works for every $L\ge2$, it identifies
the first periodic eigenmode, and it improves the quantitative obstruction
from $8/L$ to $4\sin^2(\pi/L)=O(L^{-2})$.

\begin{thebibliography}{9}

\bibitem{BalabanI}
T.~Ba\l aban,
\emph{Propagators and renormalization transformations for lattice gauge
theories. I},
Communications in Mathematical Physics \textbf{95} (1984), 17--40.
\href{https://doi.org/10.1007/BF01215753}{doi:10.1007/BF01215753}.

\bibitem{BalabanII}
T.~Ba\l aban,
\emph{Propagators and renormalization transformations for lattice gauge
theories. II},
Communications in Mathematical Physics \textbf{96} (1984), 223--250.
\href{https://doi.org/10.1007/BF01240221}{doi:10.1007/BF01240221}.

\bibitem{BalabanAveraging}
T.~Ba\l aban,
\emph{Averaging operations for lattice gauge theories},
Communications in Mathematical Physics \textbf{98} (1985), 17--51.
\href{https://doi.org/10.1007/BF01211042}{doi:10.1007/BF01211042}.

\bibitem{DimockAxial}
J.~Dimock,
\emph{Covariant axial gauge},
Letters in Mathematical Physics \textbf{105} (2015), 959--987.
\href{https://doi.org/10.1007/s11005-015-0763-0}
{doi:10.1007/s11005-015-0763-0}.

\bibitem{CombesThomas}
J.~M.~Combes and L.~Thomas,
\emph{Asymptotic behaviour of eigenfunctions for multiparticle
Schr\"odinger operators},
Communications in Mathematical Physics \textbf{34} (1973), 251--270.
\href{https://doi.org/10.1007/BF01646473}{doi:10.1007/BF01646473}.

\bibitem{Lean4}
L.~de~Moura and S.~Ullrich,
\emph{The Lean 4 theorem prover and programming language},
Automated Deduction -- CADE 28, LNCS 12699 (2021), 625--635.
\href{https://doi.org/10.1007/978-3-030-79876-5_37}
{doi:10.1007/978-3-030-79876-5\_37}.

\bibitem{mathlib}
The mathlib Community,
\emph{The Lean mathematical library},
Proceedings of CPP 2020 (2020), 367--381.
\href{https://doi.org/10.1145/3372885.3373824}
{doi:10.1145/3372885.3373824}.

\end{thebibliography}

\end{document}
